
This chapter presents the conceptual foundations for our empirical tests.
First, we recap Bybee’s \emph{Ghost in the Machine} model of belief formation in public markets,
including its formal structure and empirical implications.
Second, we explain how to adapt this framework to the illiquid, disclosure‐sparse context of private equity—without introducing a new formal model.
Finally, we develop the panel‐regression specification that uses fund‐level sentiment indices to predict future fund cash flows.

\section{Recap of Bybee’s Ghost in the Machine}

Bybee (2023) models belief formation as a two‐stage process in which agents observe textual news
and map it into forecasts via an \emph{inference machine}, capturing both rational and biased updating.

\subsection*{Information Arrival}

At each date $t$, a stream of news $N_t$ arrives. Each item $\nu\in N_t$ carries two components:
\[
m(\nu)\quad\text{(sentiment / mood)},
\quad
f(\nu)\quad\text{(fundamental signal)}.
\]

\subsection*{Inference Machine}

Analysts (and LLMs) apply a mapping $\mathcal{I}$ to aggregate past returns and textual signals into a
belief $\hat{R}_{t+1}$ about next‐period returns:
\[
  \hat{R}_{t+1}
  = \alpha\,\mathbb{E}[R_{t+1}\mid \mathcal{F}_{t-1}]
  + \beta\,f(N_t)
  + \gamma\,m(N_t)
  + \delta\,R_t
  + \varepsilon_t.
\]
Empirically, $\delta>0$ (anchoring) and $\gamma>\beta$ (sentiment > fundamentals), leading to
predictable forecast errors and mispricing in public equities.

\section{Adapting Belief Formation to Illiquid Assets}

Private‐equity (PE) differs in two key ways:

\subsection*{Sparse Disclosures}

PE funds issue lengthy, narrative reports only quarterly or semi‐annually.  Denote the set of documents at time $t$ by $D_t$; we aggregate them into a corpus
\[
C_t = \bigcup_{d\in D_t} d,
\]
analogous to $N_t$ in the public setting but arriving infrequently.

\subsection*{Long Horizons}

Realized PE outcomes $Y_{t+h}$ (e.g.\ distributions, NAV changes) materialize with lags of multiple quarters.  We therefore think of the inference machine as forecasting $\hat{Y}_{t+H}$ over horizons $H\in\{1,4\}$ quarters based on
\[
  \hat{Y}_{t+H}
  = \alpha\,\mathbb{E}[Y_{t+H}\mid\mathcal{F}_{t-1}]
  + \beta\,f(C_t)
  + \gamma\,m(C_t)
  + \delta\,Y_t
  + \varepsilon_t,
\]
where $f(C_t)$ are distilled fundamentals, $m(C_t)$ is aggregate sentiment, and $Y_t$ is last‐reported NAV or distribution.

\section{Panel Regression Model}

To test whether LLM‐extracted sentiment predicts actual PE cash flows, we estimate the following panel‐data regression:

\begin{equation}
\mathrm{Flow}_{i,t+1}
= \alpha
+ \beta\,S_{i,t}
+ \boldsymbol{X}_{i,t}'\boldsymbol{\gamma}
+ \mu_i
+ \nu_t
+ \varepsilon_{i,t+1}.
\label{eq:panel_reg}
\end{equation}

\begin{itemize}
  \item $\mathrm{Flow}_{i,t+1}$ is fund $i$’s net cash flow in quarter $t+1$ (contributions minus distributions).
  \item $S_{i,t} = m(C_{i,t})$ is the sentiment index for fund $i$ at time $t$, generated by the LLM‐based protocol.
  \item $\boldsymbol{X}_{i,t}$ is a vector of time‐varying controls (e.g.\ lagged NAV, market returns, interest rate).
  \item $\mu_i$ captures unobserved, time‐invariant fund heterogeneity (e.g.\ manager skill, strategy).
  \item $\nu_t$ captures economy‐wide shocks common to all funds in quarter $t$.
  \item $\varepsilon_{i,t+1}$ is the idiosyncratic error, satisfying $\mathbb{E}[\varepsilon_{i,t+1}\mid S_{i,t},\boldsymbol{X}_{i,t},\mu_i,\nu_t]=0$.
\end{itemize}

\subsection*{Fixed vs.\ Random Effects}

\begin{itemize}
  \item \textbf{Fixed Effects (FE):} Treat $\mu_i$ as parameters, effectively de‐mean each fund’s series.  Controls for any unobserved, time‐invariant fund traits correlated with sentiment.
  \item \textbf{Random Effects (RE):} Model $\mu_i$ as random draws, requiring $\mu_i\perp S_{i,t},\boldsymbol{X}_{i,t}$.  A Hausman test decides between FE and RE.
\end{itemize}

\subsection*{Interpretation of Coefficients}

Under the FE specification, $\beta$ measures how deviations in sentiment from a fund’s average affect its subsequent cash flows.  A positive $\beta$ implies that when sentiment is unusually high, the fund experiences greater net cash inflows (or fewer distributions) in the next quarter, consistent with optimistic beliefs driving allocation decisions.

\subsection*{Assumptions and Inference}

- We cluster standard errors by fund to allow serial autocorrelation in $\varepsilon_{i,t}$.
- We include lagged controls in $\boldsymbol{X}_{i,t}$ (e.g.\ lagged NAV, market return) to mitigate omitted‐variable bias.
- Time fixed effects $\nu_t$ absorb aggregate sentiment shocks and macro‐cycles.

\subsection{Economic Interpretation of Coefficients}

Estimating the panel regression
\[
  \mathrm{Flow}_{i,t+1}
  = \alpha
  + \beta\,S_{i,t}
  + \gamma_1\,\mathrm{Flow}_{i,t}
  + \gamma_2\,\mathrm{Rm}_{t}
  + \gamma_3\,\mathrm{Rf}_{t}
  + \mu_i
  + \nu_t
  + \varepsilon_{i,t+1}
\]
yields the following economic interpretations:

\begin{itemize}
  \item \(\beta\) (\textbf{Sentiment Effect}):
    A positive \(\beta\) implies that when the LLM‐derived sentiment index \(S_{i,t}\) is higher than its fund‐specific average, the fund experiences greater net contributions (or fewer distributions) in the next quarter.
    Economically, this suggests that optimistic narratives—captured by the sentiment index—drive LPs (or GPs) to allocate more capital to the fund, consistent with sentiment guiding real funding decisions beyond fundamentals.

  \item \(\gamma_1\) (\textbf{Flow Persistence}):
    If \(\gamma_1>0\), fund cash flows display momentum: quarters with high inflows (or low distributions) tend to be followed by similar patterns.  This reflects lifecycle effects: funds in their “investment phase” call capital more aggressively, while funds in a “harvest phase” distribute more.  A negative \(\gamma_1\) would indicate mean‐reversion in flows.

  \item \(\gamma_2\) (\textbf{Equity Market Return}):
    A positive \(\gamma_2\) indicates that during quarters when the S\&P 500 posts strong returns, PE funds also experience higher net inflows.  This captures the “wealth effect”: rising public markets increase LP wealth and risk appetite, boosting PE commitments.  A negative \(\gamma_2\) might arise if investors rotate out of PE when public equity returns are high.

  \item \(\gamma_3\) (\textbf{Risk‐Free Rate}):
    A negative \(\gamma_3\) means that a higher 10-year Treasury yield (higher cost of capital) dampens PE capital calls.  This is consistent with tighter financing conditions making investors more cautious.  Conversely, a positive \(\gamma_3\) could reflect investors chasing yield by committing more to alternatives when public‐bond yields rise.
\end{itemize}

\paragraph{Magnitude and Economic Significance}
To assess economic significance, we report “one‐standard‐deviation” impacts:
\[
  \beta \times \mathrm{SD}(S_{i,t})
  \quad\text{and}\quad
  \gamma_j \times \mathrm{SD}(X_{j,t}),
\]
which measure the change in net flows (in USD millions) associated with a one‐SD shock in each regressor.  Comparing these across regressors shows whether sentiment is a quantitatively important driver of fund flows—relative to market conditions or flow inertia.


\section{Belief-Elicitation Protocol}
Explain the pipeline: document summarization → prompt → JSON belief output.
Describe how forecasts are structured (increase, stable, decrease, etc.).
Include example prompts and schema.

\section{Construction of the Sentiment Index}
Detail how document-level forecasts are aggregated at the fund level (LP and GP sentiment indices).
Include weighting, normalization, and temporal aggregation (quarterly, rolling averages).

\section{Econometric Framework}
Present the regression specification linking sentiment to hard data (NAVs, distributions).
Include the equation, variables, and identification strategy.
Keep it conceptual here, implementation details go in next chapter.

\subsection{Panel Regression with Fund Fixed and Random Effects}

To assess whether the LLM‐generated sentiment index predicts subsequent fund cash flows, we employ a linear panel‐data regression of the form

\begin{equation}
\mathrm{Flow}_{i,t+1}
= \alpha
+ \beta\,S_{i,t}
+ \mathbf{X}_{i,t}'\boldsymbol{\gamma}
+ \mu_i
+ \nu_t
+ \varepsilon_{i,t+1},
\label{eq:panel_reg_2}
\end{equation}

where:
\begin{itemize}
  \item $\mathrm{Flow}_{i,t+1}$ is the net cash flow (or any other fund‐level outcome) for fund $i$ in period $t+1$.
  \item $S_{i,t}$ is the sentiment index for fund $i$ in period $t$.
  \item $\mathbf{X}_{i,t}$ is a vector of time‐varying control variables (see next section).
  \item $\mu_i$ is a fund‐specific effect capturing all time‐invariant heterogeneity (e.g.\ vintage, strategy, geography).
  \item $\nu_t$ is a time‐fixed effect common to all funds in period $t$ (e.g.\ aggregate market shocks).
  \item $\varepsilon_{i,t+1}$ is the idiosyncratic error term, assumed $\mathbb{E}[\varepsilon_{i,t+1}\mid S_{i,t},\mathbf{X}_{i,t},\mu_i,\nu_t]=0$.
\end{itemize}

\paragraph{Fixed Effects (FE) vs.\ Random Effects (RE)}
\begin{itemize}
  \item \textbf{FE estimator} treats $\mu_i$ as unknown parameters, effectively de‐meaning each fund’s time series and so estimating within‐fund variation.  This controls for any unobserved, time‐invariant characteristics (e.g.\ fund manager skill, mandate) that might correlate with sentiment.
  \item \textbf{RE estimator} treats $\mu_i$ as random, assuming $\mu_i\perp S_{i,t}, \mathbf{X}_{i,t}$.  Hausman’s specification test can guide the choice: if $\mu_i$ correlate with regressors, use FE; otherwise RE is efficient.
\end{itemize}

\paragraph{Identification and Interpretation}
Under the FE specification, $\beta$ is identified from the quarter–to–quarter co‐movement of sentiment and fund outcomes within each fund.  A statistically significant $\beta$ implies that departures in sentiment—above or below a fund’s own long‐run average—predictably affect its future cash flows.

\paragraph{Assumptions}
The key identifying assumption is strict exogeneity:
\[
\mathbb{E}\bigl[\varepsilon_{i,t+1}\mid S_{i,1},\dots,S_{i,t},\mu_i,\nu_{1},\dots,\nu_{t}\bigr] = 0.
\]
In practice, we will include lagged controls in $\mathbf{X}_{i,t}$ and consider clustering standard errors at the fund level to allow for serial correlation in $\varepsilon_{i,t}$.

\subsection{Control Variables}

In addition to the sentiment index \(S_{i,t}\), we include three easy‐to‐implement controls in \(\boldsymbol{X}_{i,t}\):

\begin{enumerate}
  \item \textbf{Lagged Net Cash Flow, \(\mathrm{Flow}_{i,t}\).}  
    \(\quad\)Controls for autocorrelation in each fund’s cash flows.  By including \(\mathrm{Flow}_{i,t}\), we capture momentum or mean‐reversion effects in contributions/distributions.

  \item \textbf{Market Return, \(\mathrm{Rm}_{t}\).}  
    \(\quad\)Quarterly total return on the S\&P 500 index.  This proxy for broad equity‐market conditions controls for systematic shocks that affect all funds simultaneously.

  \item \textbf{Risk‐Free Rate, \(\mathrm{Rf}_{t}\).}  
    \(\quad\)Yield on a 10-year U.S.\ zero-coupon Treasury note.  This captures the prevailing cost of capital and discount‐rate movements common to all funds.
\end{enumerate}

Thus, the full specification becomes
\[
  \mathrm{Flow}_{i,t+1}
  = \alpha
  + \beta\,S_{i,t}
  + \gamma_1\,\mathrm{Flow}_{i,t}
  + \gamma_2\,\mathrm{Rm}_{t}
  + \gamma_3\,\mathrm{Rf}_{t}
  + \mu_i
  + \nu_t
  + \varepsilon_{i,t+1}.
\]

\paragraph{Economic Rationale}  
\begin{itemize}
  \item \emph{Lagged Flows} control for predictable fund‐specific dynamics (e.g.\ vintage cycle, harvest phase).  
  \item \emph{Market Return} accounts for periods of equity‐market stress or euphoria that drive both public‐market performance and LP appetite for private deals.  
  \item \emph{Risk-Free Rate} controls for shifts in the baseline discount rate, affecting the opportunity cost of capital and general funding availability.
\end{itemize}

These three controls strike a balance between parsimony and addressing the most salient confounders.  All are available from public sources (e.g.\ Bloomberg, FRED) on a quarterly basis, and align with standard practice in panel regressions of fund‐level flows.  


\section{Validation Strategy}
Define the metrics to test predictive power and bias:
(i) forecast accuracy (RMSE, MAE),
(ii) bias coefficient (overreaction vs. fundamentals),
(iii) incremental explanatory power vs. controls.


\section{Generation of the Sentiment Panel}
Explain the aggregation results at the panel level (fund × time).
Mention the resulting dataset size, coverage, and any data cleaning.

\section{Regression Execution and Tests}
Describe how the econometric models were run (software, packages), sample periods, and variable transformations.

\section{Robustness Checks}
Discuss alternative specifications, alternative prompt settings, and sensitivity to model versions.
